\section{Introduction}
\label{sec:introduction}

\begin{figure}[!h]
\centering
\resizebox{\linewidth}{!}{%
\begin{tikzpicture}[
    x=1cm,
    y=1cm,
    font=\sffamily\scriptsize,
    entry/.style={draw=framegray, rounded corners=1.5pt, minimum width=0.62cm, minimum height=0.47cm, inner sep=1pt, fill=white},
    alias/.style={entry, draw=frameorange, very thick, fill=frameorange!10},
    panel/.style={draw=framegray!55, rounded corners=4pt, fill=framegray!3},
    flow/.style={-{Latex[length=1.7mm]}, thick, draw=framegray},
    warn/.style={draw=frameorange, rounded corners=2pt, fill=frameorange!8, align=center, inner sep=3pt}
]
% Panels
\draw[panel] (-7.85,0.72) rectangle (-0.23,-3.18);
\draw[panel] (0.23,0.72) rectangle (7.85,-3.18);
\node[font=\sffamily\small\bfseries, text=frameblue] at (-4.04,0.42) {Clean registry};
\node[font=\sffamily\small\bfseries, text=frameorange] at (4.04,0.42) {Quality-preserving refinement};

% Clean registry
\node[anchor=east] at (-6.95,-0.05) {entries};
\node[entry] (ca) at (-6.45,-0.05) {A};
\node[entry] (cb) at (-5.55,-0.05) {B};
\node[entry] (cc) at (-4.65,-0.05) {C};
\node[entry] (cd) at (-3.75,-0.05) {D};
\node[align=left, text=framegray] at (-2.00,-0.05) {$4$ entries\\$4$ roots};
\node[warn, draw=frameblue, fill=frameblue!7, text=black, minimum width=2.65cm] (cleanmass) at (-5.75,-1.12) {per-entry mass\\$=1/4$ per root};
\node[warn, draw=frameblue, fill=frameblue!7, text=black, minimum width=2.65cm] (cleanpath) at (-2.25,-1.12) {query transcript\\$x_2\!\rightarrow x_5\!\rightarrow x_9$};
\draw[flow] (cleanmass) -- (cleanpath);
\node[font=\bfseries, text=framegreen] at (-4.04,-2.22) {lower cumulative regret};
\node[text=framegray, align=center] at (-4.04,-2.75) {candidate quality vectors unchanged};

% Refined registry
\node[anchor=east] at (1.20,-0.05) {entries};
\node[alias] (ra) at (1.52,-0.05) {A};
\node[alias] (ra1) at (2.25,-0.05) {$A_1$};
\node[alias] (ra2) at (2.98,-0.05) {$A_2$};
\node[alias] (ra3) at (3.71,-0.05) {$A_3$};
\node[alias] (ra4) at (4.44,-0.05) {$A_4$};
\node[entry] (rb) at (5.25,-0.05) {B};
\node[entry] (rc) at (5.98,-0.05) {C};
\node[entry] (rd) at (6.71,-0.05) {D};
\node[draw=frameorange, dashed, rounded corners=2pt, fit=(ra)(ra1)(ra2)(ra3)(ra4), inner sep=2pt,
      label={[text=frameorange, font=\sffamily\tiny]below:one accountable root}] {};
\node[warn, minimum width=2.80cm] (refmass) at (2.20,-1.40) {per-entry mass\\A-family $=5/8$};
\node[warn, minimum width=2.80cm] (refpath) at (5.72,-1.40) {different transcript\\$x_1\!\rightarrow x_4\!\rightarrow x_8$};
\draw[flow, draw=frameorange] (refmass) -- (refpath);
\node[font=\bfseries, text=frameorange] at (4.04,-2.35) {higher cumulative regret can result};
\node[text=framegray, align=center] at (4.04,-2.83) {same roots and per-example utility};

% Missing observation boundary
\node[warn, minimum width=12.5cm, font=\sffamily\scriptsize] at (0,-3.78)
{Predictions expose submitted response rows, but not the accountable map $\rho:E\!\rightarrow\!R$ or which entries deserve independent evidential mass.};
\end{tikzpicture}
}%
\caption{The candidate list is an empirical prior. Adding entries from one accountable root can change active queries even when the added entries introduce no new per-example quality vector. The observable responses do not authenticate the root map. This schematic illustrates the mechanism; Tables~\ref{tab:primary-llm} and~\ref{tab:step88} report measured effects.}
\label{fig:registry-refinement}
\end{figure}


Expert labels are costly, so DiffUse, MODEL SELECTOR, CODA, LLM Selector, and Select-LLM use candidate outputs, weak judges, or consensus to prioritize labels for model comparison \citep{ashurytahan2024diffuse,okanovic2025modelselector,kay2025coda,durmazkeser2025llmselector,durmazkeser2026selectllm}. In the candidate-indexed Bayesian and consensus mechanisms we audit, however, every listed entry receives separate probability or consensus mass. They state no invariance requirement for two lists encoding the same behavioral alternatives.

The missing object is the \emph{candidate evidence frame}: a map from listed \emph{submitted entries} to accountable \emph{candidate roots}, together with a root-mass rule. Predictions alone need not reveal the ownership or lineage evidence needed to authenticate this frame. Consequently, treating entries as the population makes registry representation part of the experiment.

The issue occurs naturally. In our frozen roster, \texttt{claude-instant-1.2} and \texttt{claude-instant-v1} have identical responses and scores on all 2,500 task items. Restoring the latter to its 30-entry quotient raises one behavior's mass from $1/30$ to $2/31$ and changes every one of 1,500 Select-LLM paths; random querying is unchanged. Because unique-label substitution is only 8.7/4.7/3.2\% of budget and terminal harm is absent, this is evidence of representation dependence, not broad natural damage.

We organize the paper around three questions. \emph{Is the frame necessary?} We prove a tight budget-linear cost for policies that cannot observe it and construct four aliases that redirect every query in a Select-LLM instance. \emph{Does refinement causally change model-selection cost?} Public-reference T2 aliases preserve correctness coordinate-wise while all existing candidates, labels, pools, and code remain fixed. They increase cumulative regret on three tasks and terminal regret by 2.91/10.56 points on MedQA/GSM8K; a fixed-query counterfactual removes the effect. \emph{Can the channel operate without target outcomes?} Calibration-only variants pass three of eight untouched scenarios, and raw-text learned adapters yield $+.702$ [$+.615,+.788$] terminal points on a disjoint 13,000-row IMDB outcome. A stricter one-shot Yelp transfer is positive but below its predeclared materiality threshold. Cross-selector and untargeted audits show selective rather than generic degradation.

Sequential multiplicity is known in strategic-replication bandits, perturbed-clone active evaluation, redundant-committee quotients, and static leaderboards \citep{shin2022strategicreplication,shin2025replicationproof,lanctot2026activeevaluation,nguyen2025unreal,singh2025leaderboard,hays2026strategic}. We claim neither the first sequential clone effect nor the first multiplicity attack. Our remaining object is specific: in shared-label model selection with a latent root map, entry mass changes which reusable references are acquired and hence root-level regret. A static rank can be recomputed after quotienting; an acquisition transcript cannot recover labels never requested. Exact hashes remove literal copies, while response similarity cannot authenticate ownership.

Our contributions are:
\begin{enumerate}
    \item We define evidence-frame ambiguity and show both its information-theoretic cost and a constant-size, full-path amplification.
    \item Exact-quality interventions and fixed-query counterfactuals establish the complete refinement $\rightarrow$ acquisition $\rightarrow$ evidence $\rightarrow$ root-cost chain; a natural alias anchors the representation problem in an existing roster.
    \item Reference-free, learned, cross-selector, defense, and untargeted audits identify where the effect transfers and where it does not. They support a targeted evaluation-integrity claim, not universal or production harm.
\end{enumerate}

We do not claim that every active selector is vulnerable, that a natural alias is adversarial, that every refinement is harmful, or that similarity reveals ownership. The resulting reliability requirement is narrower: an active evaluation should either bind submitted entries to a declared candidate evidence frame or qualify conclusions that depend on unaudited entry multiplicity.
